\clearpage
\fancyhead[L]{\small SCHWINGUNGSDAUER EINER MAGNETNADEL.}
\SourceLeaf{380}
\pdfbookmark[0]{Anleitung zur Bestimmung der Schwingungsdauer einer Magnetnadel. 1837}{schwingungsdauer-1837}
\begin{center}
{\large ANLEITUNG ZUR BESTIMMUNG DER SCHWINGUNGSDAUER EINER MAGNETNADEL.}
\par\medskip
Resultate aus den Beobachtungen des magnetischen Vereins. 1837. IV.
\end{center}

Die Aufgabe, zu deren Lösung hier eine Anleitung gegeben werden soll, ist in mehrfacher Hinsicht von Interesse. Eine gewisse, wenn auch noch nicht sehr genaue Kenntnis der Schwingungsdauer braucht man bereits, um die Anweisungen zur Bestimmung der Ruhelage der Nadel anzuwenden (\textit{Res.} von 1836, II.). Um dagegen die absolute Intensität des Erdmagnetismus zu ermitteln, ist der mit größtmöglicher Genauigkeit bestimmte Wert der Schwingungsdauer ein wesentliches Element. Aber auch für sich genommen kann die Ausführung der Arbeitsschritte zur Bestimmung der Schwingungsdauer als gute Vorbereitung auf astronomische Beobachtungen gelten. Diese Schritte haben besonders große Ähnlichkeit mit den Beobachtungen von Sterndurchgängen am Meridianfernrohr. Sie bieten den Vorteil, dass sie genauer ausgeführt werden können und sich, von ungünstigen Luftverhältnissen ungestört, zu jeder beliebigen Stunde vornehmen lassen. Solche Beobachtungen am Magnetometer scheinen daher auch besonders geeignet, einen noch nicht hinreichend aufgeklärten Gegenstand zu erhellen. Die Ergebnisse verschiedener Beobachter am Meridianfernrohr unterscheiden sich um konstante Beträge. Das kann nur daran liegen, dass die optischen oder akustischen Eindrücke oder beide bei verschiedenen Personen und je nach den Umständen nicht gleichzeitig ins Bewusstsein gelangen.

Die Schwingungsdauer einer Nadel ist die Zeit zwischen zwei aufeinanderfolgenden äußersten Stellungen (Elongationen). Streng genommen \SourceLeaf{381}verweilt die Nadel in keiner Elongation. Da die Bewegungsgeschwindigkeit aber stetig abnimmt, bis sie verschwindet, und von da an ebenso stetig wieder zunimmt, scheint die Nadel für unsere Sinne um den Zeitpunkt der Elongation herum mehr oder weniger lange stillzustehen. Bei kurzen Schwingungszeiten und großen Bögen ist diese scheinbare Ruhezeit freilich kaum als solche erkennbar. In jedem Fall bleibt eine solche unmittelbare Wahrnehmung des Zeitpunkts der Elongation an Genauigkeit weit hinter einer mittelbaren Bestimmung durch einander entsprechende Beobachtungen zurück. Man bildet nämlich den Mittelwert der beiden Zeitpunkte, an denen ein und derselbe Teilstrich der Skala beim Hin- und Rückgang auf dem senkrechten Faden des Fernrohrs erscheint.

Im Allgemeinen wählt man dafür am besten einen Teilstrich in oder nahe der Mitte des Schwingungsbogens. Dort ist die Bewegung am schnellsten, sodass sich der Zeitpunkt selbst am genauesten beobachten lässt. Außerdem werden bei der Beobachtung mehrerer aufeinanderfolgender Schwingungen die Zeitabstände zwischen den Aufzeichnungen dann nahezu gleich. In einzelnen Fällen, namentlich bei sehr langer Schwingungsdauer, kann es allerdings gelegentlich vorteilhaft sein, eine andere oder sogar mehrere verschiedene Stellen der Skala zu verwenden. Das kann hier jedoch außer Betracht bleiben.

Bei kleinen Schwingungen wählt man am besten den Teilstrich, der der Mitte am nächsten liegt. Bei größeren Schwingungen sollte man den leichter zu beobachtenden nächstgelegenen Fünfer- oder Zehnerstrich der Skala vorziehen. Bei sehr großen Schwingungen geht die Mitte der Skala jedoch so schnell durch das Gesichtsfeld, dass man die gewählte Stelle auf der Skala hinreichend auffällig machen muss, etwa durch einen nicht zu feinen schwarzen Faden, den man über die Skala hängt.

Fällt der Durchgang am Faden nicht genau mit einem Sekundenschlag zusammen, so schätzt man den Sekundenbruchteil nach dem Verhältnis der beiden Abstände, in denen die betreffende Stelle der Skala beim vorhergehenden und beim folgenden Sekundenschlag vom Faden erscheint. So sind es auch die meisten Astronomen bei der Beobachtung am Meridianfernrohr gewohnt. Unmittelbar eingeteilt wird also der Raum. Die Zeit wird dabei nur mittelbar eingeteilt.

Die Schwingungsdauer lässt sich aus einer einzigen Schwingung natürlich nur mit sehr begrenzter Genauigkeit bestimmen. Man stützt ihre Bestimmung daher immer auf eine größere Zahl aufeinanderfolgender Schwingungen. Zwar hängt
\expandafter\def\csname HistoricalNoteText391\endcsname{Sind die Teile des schwingenden Apparats nicht fest miteinander verbunden, so nehmen die Schwingungsbögen immer rasch und unregelmäßig ab.}
\expandafter\def\csname HistoricalNoteText382\endcsname{Es braucht kaum erwähnt zu werden, dass bei rasch wechselnden starken Änderungen der Deklination \textit{sehr kleine} Schwingungen (die sich schon für sich genommen wenig zur Bestimmung der Schwingungsdauer eignen) eine ganz entstellte Dauer zeigen können.}
\SourceLeaf{382}die Schwingungsdauer von veränderlichen Elementen ab und ändert sich daher, streng genommen, selbst beständig. Von ganz ungewöhnlichen Fällen abgesehen\HistoricalNote{382}, kann man diese Veränderlichkeit über einen nicht allzu kurzen Zeitraum als völlig unmerklich ansehen. Jedenfalls entspricht die aus einer beträchtlichen Zahl von Schwingungen abgeleitete Dauer einer einzelnen Schwingung den Mittelwerten der einzelnen in Betracht kommenden Elemente während dieser Zeit. Es ist aber offensichtlich überhaupt nicht nötig, den Bewegungen der Nadel während eines solchen Zeitraums ununterbrochen zu folgen. Es genügt, die Zeitpunkte der ersten und der letzten Elongation zu kennen, sobald die Schwingungsdauer näherungsweise so genau bekannt ist, dass über die \textit{Anzahl} der Schwingungen im verstrichenen Zeitraum kein Zweifel mehr bestehen kann. Das wird noch dadurch erleichtert, dass man immer im Voraus weiß, ob diese Anzahl gerade oder ungerade ist. Das ergibt sich daraus, ob die erste und die letzte Elongation gleichnamig oder ungleichnamig sind.

Wenn die Schwingungsdauer nicht allzu kurz ist, können zwischen den Durchgängen auch die Elongationen selbst (nämlich die äußersten Ablesewerte auf der Skala) mit aufgezeichnet werden. So lassen sich daraus die Amplituden ableiten, die man für eine genauere Berechnung der Schwingungsdauer braucht. Ihre fortschreitende Abnahme gibt außerdem für sich genommen Anlass zu bemerkenswerten Überlegungen.

Wie man die Beobachtungen auswertet, um Ergebnisse daraus zu gewinnen, lässt sich am besten an einem Beispiel zeigen. Dafür werden hier Beobachtungen vom 29. November 1835 am Magnetometer der Sternwarte mit fünfundzwanzig Pfund schwerem Stab gewählt. Die folgende Tafel I. enthält zunächst die Rohbeobachtungen.

\begin{center}
\begingroup\fontsize{9}{11}\selectfont\setlength{\tabcolsep}{2pt}
\begin{tabular}{r|r||r|r||r|r||r|r}
& $1755.1$ && $497.8$ && $645.9$ && $1232.1$ \\[-4pt]
$21^{\mathrm{h}}55'26''9$ && $23^{\mathrm{h}}38'49''2$ && $1^{\mathrm{h}}10'12''6$ && $2^{\mathrm{h}}49'19''7$ & \\[-4pt]
& $266.0$ && $1502.2$ && $1341.5$ && $775.9$ \\[-4pt]
$56\quad8.4$ && $39\quad31.5$ && $54.2$ && $50\quad1.5$ & \\[-4pt]
& $1751.8$ && $500.1$ && $647.3$ && $1231.0$ \\[-4pt]
$51.2$ && $40\quad13.6$ && $11\quad37.0$ && $44.1$ & \\[-4pt]
& $268.5$ && $1499.1$ && $1339.4$ && $776.4$ \\[-4pt]
$57\quad33.0$ && $56.0$ && $12\quad18.4$ && $51\quad25.8$ & \\[-4pt]
& $1748.9$ && $502.6$ && $648.7$ && $1228.7$ \\[-4pt]
$58\quad15.5$ && $41\quad38.1$ && $13\quad1.3$ && $52\quad8.5$ & \\[-4pt]
& $271.6$ && $1496.5$ && $1337.0$ && $778.0$ \\[-4pt]
$57.4$ && $42\quad20.3$ && $43.0$ && $50.0$ & \\[-4pt]
& $1744.2$ && $506.0$ && $650.7$ && $1227.0$
\end{tabular}\endgroup
\end{center}

\SourceLeaf{383}Diese Beobachtungen bestehen, wie man sieht, aus vier Sätzen. Die erste Spalte enthält die Zeitpunkte, zu denen der Skalenpunkt 1000 am Faden vorbeiging, die zweite die Ablesewerte der Elongationen. Diesmal beginnen die Werte mit der Elongation vor dem ersten Durchgang und enden mit der Elongation nach dem letzten Durchgang. Beobachtet man die Elongationen nicht mit, so vermerkt man am besten bei jedem Durchgang, ob zunehmende oder abnehmende Zahlen am Faden vorbeigingen. In Göttingen geht man dabei wie hier gezeigt vor.
\begin{center}
\begin{tabular}{l}
$21^{\mathrm{h}}55'26''9-$ \\
$\phantom{21^{\mathrm{h}}}56\quad8.4+$ \\
und so fort \\
$23^{\mathrm{h}}38'49''2+$ \\
und so weiter
\end{tabular}
\end{center}

Das ist nötig, um unterscheiden zu können, welche der aus den Durchgängen abgeleiteten Elongationszeiten zu Minima und welche zu Maxima gehören. Bei der Nummerierung der Elongationszeiten haben wir es uns zur Gewohnheit gemacht, die ersteren mit geraden und die letzteren mit ungeraden Zahlen zu bezeichnen. Vor die aus den beiden ersten Durchgängen abgeleitete Elongationszeit $21^{\mathrm{h}}55'47''65$ setzt man daher die Zahl 0 und so fort.

Die folgende Tafel II. enthält nun die nächsten aus den unmittelbaren Beobachtungen berechneten Ergebnisse.
\begin{center}
\begingroup\fontsize{9}{11}\selectfont\setlength{\tabcolsep}{4pt}
\begin{tabular}{r|r|r||r|r|r}
0 & $21^{\mathrm{h}}55'47''65$ & $21^{\mathrm{h}}55'47''65$ & 277 & $1^{\mathrm{h}}10'33''40$ & $21^{\mathrm{h}}55'49''54$ \\
1 & $56\quad29.80$ & $47.62$ & 278 & $11\quad15.60$ & $49.56$ \\
2 & $57\quad12.10$ & $47.74$ & 279 & $11\quad57.70$ & $49.48$ \\
3 & $57\quad54.25$ & $47.71$ & 280 & $12\quad39.85$ & $49.45$ \\
4 & $58\quad36.45$ & $47.73$ & 281 & $13\quad22.15$ & $49.57$ \\[7pt]
147 & $23\quad39\quad10.35$ & $21\quad55\quad49.89$ & 418 & $2\quad49\quad40.60$ & $21\quad55\quad49.36$ \\
148 & $39\quad52.55$ & $49.91$ & 419 & $50\quad22.80$ & $49.38$ \\
149 & $40\quad34.80$ & $49.98$ & 420 & $51\quad4.95$ & $49.35$ \\
150 & $41\quad17.05$ & $50.05$ & 421 & $51\quad47.15$ & $49.37$ \\
151 & $41\quad59.20$ & $50.02$ & 422 & $52\quad29.25$ & $49.29$
\end{tabular}\endgroup
\end{center}

Die zweite Spalte enthält hier die sich ergebenden Elongationszeiten. Die Nummerierung in der ersten Spalte ergibt sich für die ersten fünf Elongationen von selbst. Für die späteren findet man sie auf folgende Weise.

\SourceLeaf{384}Der Vergleich der. Elongation 0 mit 4 ergibt als Näherungswert der Schwingungsdauer $42''20$. Teilt man die Zeit zwischen der Elongation 4 und der nächsten, $1^{\mathrm{h}}40'33''90$, durch diesen Wert und bedenkt, dass die Ordnungszahl der nächsten Elongation ungerade sein muss, so lässt der Quotient $142.983$ keinen Zweifel daran, dass zwischen den beiden Elongationen 143 Schwingungen verstrichen sein müssen. Nähme man nämlich 141 oder 145 an, ergäbe sich eine Schwingungsdauer von $42''7936$ beziehungsweise $41''6131$. Beide weichen zu stark vom Näherungswert $42''20$ ab, um zulässig zu sein. Ausgehend von 143 Schwingungen findet man die Schwingungsdauer $42''1951$. Diesen Wert könnte man beim Übergang zu dem folgenden Beobachtungssätzen zugrunde legen, um ihre Nummerierung zu bestimmen. Im vorliegenden Fall, in dem keine \textit{sehr} langen Unterbrechungen vorkommen, reicht allerdings schon der erste Näherungswert überall aus.

Um die Schwingungsdauer genauer zu bestimmen und sogar ihre Veränderlichkeit im Verlauf der ganzen Beobachtungsreihe zu erkennen, kann man zunächst den ersten Satz auf folgende Weise mit dem zweiten vergleichen. Für die Dauer von 147 Schwingungen erhält man
\begin{center}
\begin{tabular}{rcl}
$0\text{—}147$ & $.\ .\ .\ .\ .$ & $1^{\mathrm{h}}43'22''70$\\
$1\text{—}148$ & & $\phantom{1^{\mathrm{h}}43'}22.75$\\
$2\text{—}149$ & & $\phantom{1^{\mathrm{h}}43'}22.70$\\
$3\text{—}150$ & & $\phantom{1^{\mathrm{h}}43'}22.80$\\
$4\text{—}151$ & & $\phantom{1^{\mathrm{h}}43'}22.75$
\end{tabular}
\end{center}
und damit den Mittelwert $1^{\mathrm{h}}43'22''74$. Die Dauer einer Schwingung ist also $=42''19551$. Ebenso erhält man die Schwingungsdauer zwischen dem zweiten und dritten Satz $=42''17654$, zwischen dem dritten und vierten $=42''17879$ und zwischen dem ersten und vierten, also als Mittel der ganzen Reihe, $=42''18344$.

Man kann diese Rechnung auch in etwas abgewandelter Form führen. Das hat zugleich den Vorteil, dass sich der regelmäßige Verlauf aller einzelnen Beobachtungen klar überblicken lässt. Zunächst reduziert man die einzelnen gefundenen Elongationszeiten mit einem Näherungswert der Schwingungsdauer auf dieselbe Epoche. Dazu rechnet man von jeder dieser Zeiten mithilfe des Näherungswerts um die gesamte Dauer der seit dieser Epoche verstrichenen Schwingungen zurück. Man zieht also von jeder Zahl der zweiten
\SourceLeaf{385}Spalte das Produkt aus dem angenommenen Wert und der entsprechenden Zahl der ersten Spalte ab. Wären die Beobachtungen absolut genau und wäre die Schwingungsdauer genau konstant und genau gleich dem angenommenen Wert, müssten sämtliche so reduzierten Zahlen genau gleich ausfallen. Steigen die Zahlen von einem Satz zum nächsten, so erkennt man daran, dass die zugrunde gelegte Schwingungsdauer für diesen Zeitraum zu klein war, und umgekehrt. Das unregelmäßige Schwanken der Ergebnisse innerhalb ein und desselben Satzes bietet dagegen einen Maßstab für die Genauigkeit der Beobachtungen selbst.

In unserem Beispiel ergibt der Vergleich der ersten mit der letzten Elongationszeit die Schwingungsdauer $=42''18389$. Statt dieses Werts wurde bei der Berechnung der Zahlen der dritten Spalte der Näherungswert $42''18$ zugrunde gelegt. Man sieht so auf einen Blick, dass diese Schwingungsdauer für die Zeit vom ersten zum zweiten Satz etwas zu klein ist, dagegen vom zweiten zum dritten. und ebenso vom dritten zum vierten ein wenig zu groß. Um genaue Ergebnisse zu erhalten, bildet man den Mittelwert der Zahlen der dritten Spalte für jeden Satz. Diese Mittelwerte
\begin{center}
\begin{tabular}{r}
$21^{\mathrm{h}}55'47''69$\\
$49.97$\\
$49.52$\\
$49.35$
\end{tabular}
\end{center}
kann man als genauere Werte der reduzierten Zeiten zu den Ordnungszahlen 2, 149, 279, 420 ansehen. Die Beobachtungen haben also vom ersten zum zweiten Satz einen Vorsprung von $2''28$ gegenüber dem für 147 Schwingungen angenommenen Verlauf. Das ergibt für eine Schwingung $0''01551$, sodass der korrigierte Wert $42''19551$ genau mit dem oben gefundenen übereinstimmt. Der Vergleich der folgenden Sätze ergibt dasselbe Resultat.

Schon ein Blick auf die Zahlen, die zu ein und demselben Satz gehören, zeigt die Güte der Beobachtungen. Dennoch seien hier die Regeln angegeben, nach denen sich in geeigneten Fällen der Maßstab für die Genauigkeit genauer ermitteln lässt. Bezeichnen wir den mittleren bei einem Durchgang zu erwartenden Fehler mit $\varepsilon$, die Anzahl der Ergebnisse eines Satzes mit $s$ und die Summe der Quadrate der Abweichungen dieser einzelnen Ergebnisse
\SourceLeaf{386}von ihrem Mittelwert mit $q$, so können wir näherungsweise annehmen
\[
\frac{(s-1)^2\varepsilon\varepsilon}{2s}=q
\]
oder, wenn mehrere Sätze vorliegen,
\[
\varepsilon\varepsilon\sum\frac{(s-1)^2}{2s}=\sum q
\]
also
\[
\varepsilon=\BareRoot{\frac{\sum q}{\sum\frac{(s-1)^2}{2s}}}
\]

In unserem Beispiel betragen die Abweichungen vom Mittelwert beim ersten Satz $0''04$, $0''07$, $0''05$, $0''02$, $0''04$, also ist $q=110$, wenn man eine Hundertstelsekunde als Einheit nimmt. Ferner ist $s=5$, also $\frac{(s-1)^2}{2s}=\frac85$. Für die drei folgenden Sätze gilt beim gleichen Wert von $s$, dass $q=190;\ 110;\ 50$. Wir haben also
\[
\tfrac{32}{5}\varepsilon\varepsilon=460
\]
oder
\[
\varepsilon=8.5,\qquad\text{d. i. }\varepsilon=0''085.
\]

Allerdings ist zu beachten, dass die Gültigkeit dieser Regel von mehreren Bedingungen abhängt, die unser Beispiel nicht hinreichend erfüllt. Erstens muss sich der angenommene Näherungswert der Schwingungsdauer, mit dem die reduzierten Zahlen berechnet wurden, ohne merklichen Fehler für jeden Satz als der wahre Wert ansehen lassen. Zweitens müssen die verschiedenen Sätze, die man zusammenfasst, unter annähernd gleichen Umständen beobachtet worden sein, soweit diese die Genauigkeit der Beobachtung beeinflussen können. Beides trifft in unserem Beispiel nicht zu, und die obige Rechnung ist daher nur als Erläuterung der Formel anzusehen. Will man genauere Bestimmungen, so stellt man am besten zunächst eigens zu diesem Zweck Beobachtungen an. Ohne einer späteren ausführlicheren Behandlung dieses Gegenstands vorzugreifen, sei hier nur bemerkt, dass die Genauigkeit der Beobachtung außer von den Eigenheiten des Apparats und des Beobachters auch von der besseren oder schlechteren Beleuchtung der Skala, der Geschwindigkeit der Schwingungsbewegung und der Beschaffenheit der Uhr abhängt. Sowohl eine allzu schnelle als auch eine allzu langsame Bewegung ist für die Genauigkeit des Be\SourceLeaf{387}obachtens weniger günstig als eine mittlere Geschwindigkeit. Mit einer Sekundenuhr beobachtet man weniger genau als mit einem Chronometer, das kürzere Zeitabschnitte durch Schläge angibt. Unter den günstigsten Umständen übertrifft die Genauigkeit dieser Beobachtungen die der besten Beobachtungen am Meridianfernrohr bei Weitem.

Unter sonst gleichen Umständen ist die Schwingungsdauer bekanntlich umso kürzer, je kleiner der Schwingungsbogen ist. Nähert sich der Bogen immer weiter der Null, so nähert sich die Dauer einem Grenzwert. Bezeichnen wir diesen Grenzwert, oder nach der üblichen Ausdrucksweise die Dauer einer unendlich kleinen Schwingung, mit $T$ und die einem Schwingungsbogen $G$ entsprechende Dauer mit $T'$, so gilt bekanntlich
\[
T'=T(1+\tfrac14\sin\tfrac14G^2+\tfrac14.\tfrac9{16}\sin\tfrac14G^4+\tfrac14.\tfrac9{16}.\tfrac{25}{36}\sin\tfrac14G^6+\text{etc.})
\]

Da man beim Gebrauch des Magnetometers auf kleine Bögen beschränkt ist, kann man die Glieder vierter und höherer Ordnung unbedenklich weglassen und deshalb auch $\frac{g}{8r}$ anstelle von $\sin\tfrac14G$ setzen. Dabei bezeichnet $g$ das dem Bogen $G$ entsprechende Stück der Skala und $r$ die horizontale Entfernung der Skalenmitte vom Spiegel. Wir haben also $T'=T(1+\frac{gg}{256rr})$ oder mit derselben Genauigkeit $T=T'(1-\frac{gg}{256rr})$. In unserem Beispiel ist in Skalenteilen oder Millimetern $r=4775.9$. Der Schwingungsbogen zwischen den Elongationen 0 und 1 ist $=1485.8$. Damit beträgt die Reduktion der Schwingungsdauer auf die einer unendlich kleinen Schwingung $=-0''01595$. Ebenso findet man als Reduktion der zweiten Schwingung $0''01590$, der dritten $0''01584$ und der vierten $0''01577$. Im Mittel aus den vier ersten Schwingungen ergibt sich so für die auf unendlich kleine Schwingungen reduzierte Dauer $=42''18414$.

Dieses Verfahren setzt aber voraus, dass die Elongationen ununterbrochen beobachtet werden. Bei einer Schwingungsreihe, von der nur Anfang und Ende beobachtet wurden, kann man die Reduktion auf unendlich kleine Bögen notfalls mit einem mittleren Wert für die Größe des Schwingungsbogens vornehmen, wenn der Bogen in der Zwischenzeit nur mäßig abgenommen hat. Für die Reduktion einer längeren Reihe solcher Schwingungen, in der der Bogen stark abgenommen hat, braucht man jedoch unbedingt eine wenigstens näherungsweise richtige
\SourceLeaf{388}Kenntnis des Gesetzes, nach dem diese Abnahme erfolgt. Ich gehe daher zur Auswertung der beobachteten Elongationen über, deren erste Ergebnisse die folgende Tafel III. enthält.
\begin{center}
\fontsize{9}{11}\selectfont
$\left.\begin{array}{r|r|r|r}
0&1009.725&1487.45&3.17244\\
1&1009.525&1484.55&3.17160\\
2&1009.425&1481.85&3.17081\\
3&1009.475&1478.85&3.16992\\
4&1009.075&1474.95&3.16878
\end{array}\right\}3.170710$\\[6pt]
$\left.\begin{array}{r|r|r|r}
147&1000.575&1003.25&3.00141\\
148&1000.375&1000.55&3.00024\\
149&1000.225&997.75&2.99902\\
150&1000.200&995.20&2.99791\\
151&1000.400&992.20&2.99660
\end{array}\right\}2.999036$\\[6pt]
$\left.\begin{array}{r|r|r|r}
277&994.050&694.90&2.84192\\
278&993.875&693.15&2.84083\\
279&993.700&691.40&2.83973\\
280&993.450&689.50&2.83853\\
281&993.350&687.30&2.83715
\end{array}\right\}3.839630$\\[6pt]
$\left.\begin{array}{r|r|r|r}
418&1003.725&455.65&2.65863\\
419&1003.575&454.85&2.65787\\
420&1003.125&453.45&2.65653\\
421&1002.950&451.50&2.65466\\
422&1002.925&449.85&2.65307
\end{array}\right\}2.656152$
\end{center}

Die erste Spalte enthält die Ordnungszahl jeder Elongation. Die zweite gibt die entsprechende Ruhelage der Nadel nach der Formel $\tfrac14(a+2b+c)$ an, wobei $b$ die beobachtete Elongation, $a$ und $c$ die vorhergehende und die folgende bezeichnen (vgl. \textit{Resultate} für 1836, II). Die dritte Spalte enthält den doppelten Abstand jeder Elongation von der entsprechenden Ruhelage, also den Bogen, der von dort aus ohne die Ursachen beschrieben worden wäre, die ihn zu verkleinern suchen. Das ist $\tfrac12(a-2b+c)$ oder $\tfrac12(2b-a-c)$, also das Mittel aus dem vorhergehenden und dem folgenden Schwingungsbogen. Die vierte Spalte enthält den Logarithmus dieser Zahl. Daneben steht schließlich der Mittelwert aus den Zahlen der vierten Spalte, die zu einem Satz gehören.

Alle Erfahrungen sprechen dafür, dass man wenigstens während eines mäßig langen Zeitraums die Zahlen der dritten Spalte als in geometrischer
\SourceLeaf{389}und ihre Logarithmen folglich als in arithmetischer Folge abnehmend betrachten oder dies wenigstens als plausibelste Näherung gelten lassen darf. Die kleinen Unregelmäßigkeiten beim Vergleich aufeinanderfolgender Zahlen eines Satzes sind nur den unvermeidlichen kleinen Beobachtungsfehlern oder kleinen zufälligen Störungen zuzuschreiben. Ihren nachteiligen Einfluss verringert man so weit wie möglich, indem man die Mittelwerte der fünften Spalte den entsprechenden mittleren Ordnungszahlen zuordnet und daraus den Verlauf der ganzen Beobachtungsreihe ableitet.

Wir erhalten demnach als Logarithmen der Amplituden zu den Elongationen
\begin{center}
\begin{tabular}{rcl}
$2$ & $.\ .\ .\ .\ .\ .\ .\ .\ .\ .$ & $3.170710$\\
$149$ & $.\ .\ .\ .\ .\ .\ .\ .\ .\ .$ & $2.999036$\\
$279$ & $.\ .\ .\ .\ .\ .\ .\ .\ .\ .$ & $2.839630$\\
$420$ & $.\ .\ .\ .\ .\ .\ .\ .\ .\ .$ & $2.656152$
\end{tabular}
\end{center}

Der Logarithmus hat also vom ersten zum zweiten Satz während 147 Schwingungen um $0.171674$ abgenommen. Gleichmäßig aufgeteilt, ergibt das $0.00116785$ für eine Schwingung. Diesen Quotienten nenne ich das \textit{logarithmische Dekrement}. Vom zweiten zum dritten Satz ergibt es sich $=0.00122620$, vom dritten zum vierten $=0.00130126$. Die Abnahme ist hier also zumindest annähernd gleichförmig. Auf die Veränderlichkeit des logarithmischen Dekrements werde ich weiter unten zurückkommen.

Unter der Voraussetzung, dass die Amplituden während einer Schwingungsreihe in geometrischer Folge abgenommen haben, lassen sich die Schwingungen leicht auf unendlich kleine reduzieren. Ist $g$ die Größe der ersten Schwingung in Skalenteilen, $g^0$ die Größe der letzten und $\theta$ der Exponent der geometrischen Folge, also $g^0=g\theta^\mu$, wobei $\mu$ die Anzahl der Schwingungen bezeichnet, so beträgt die Reduktion der ersten Schwingungsdauer auf die einer unendlich kleinen Schwingung
\[
-\frac{Tgg}{256rr},
\]
die der zweiten
\[
=-\frac{Tgg\theta\theta}{256rr}
\]
und so weiter. Die Summe aller Reduktionen ist also
\[
=-\frac{T(gg-g^0g^0\theta\theta)}{256rr(1-\theta\theta)}
\]

\SourceLeaf{390}Bezeichnen wir die zur Anfangselongation und zur Endelongation gehörenden Amplituden, auf dieselbe Weise berechnet wie in Tafel III, mit $h$ und $h^0$, so ist $h=\tfrac12(\frac{g}{\theta}+g)$, $h^0=\tfrac12(g^0+g^0\theta)$, mithin ist die obige Summe
\[
=-\frac{T(hh-h^0h^0)}{64rr}.\frac{\theta\theta}{(1+\theta)^2(1-\theta\theta)}
\]

In allen hier betrachteten Fällen unterscheidet sich $\theta$ nur wenig von eins. Das mit $\lambda$ bezeichnete logarithmische Dekrement $=\log\frac1\theta$ ist also ein kleiner Bruch. Für den zweiten Faktor in jenem Ausdruck ergibt sich, wenn $m$ den Modul des Logarithmensystems bezeichnet, die Reihe
\[
\frac{\theta\theta}{(1+\theta)^2(1-\theta\theta)}=\frac{m}{8\lambda}-\tfrac5{96}.\frac{\lambda}{m}+\tfrac{37}{2880}(\frac{\lambda}{m})^3\ \text{etc.}
\]
Man kann daher mit hinreichender Genauigkeit diesen Faktor durch das erste Glied $\frac{m}{8\lambda}$ allein ersetzen. Die Reduktion der Gesamtdauer der $\mu$ Schwingungen beträgt also
\[
=-\frac{Tm(hh-h^0h^0)}{512rr\lambda}
\]
und die Reduktion des Mittelwerts für eine Schwingung beträgt
\[
=-\frac{Tm(hh-h^0h^0)}{512rr\lambda\mu}
\]

Aus den oben angeführten Gründen ist es bei diesen Formeln gleichgültig, ob man für $T$ den berichtigten oder den unberichtigten Wert einsetzt. In unserem Beispiel ergeben sich die folgenden Werte für die Reduktion
\begin{center}
\fontsize{9}{11}\selectfont
\begin{tabular}{r|r|r|r}
& \multicolumn{1}{c|}{für die ganze} & \multicolumn{1}{c|}{für eine} & \multicolumn{1}{c}{Reduzierter}\\
& \multicolumn{1}{c|}{Zeit} & \multicolumn{1}{c|}{Schwingung} & \multicolumn{1}{c}{Wert}\\[5pt]
$2\ .\ .\ .\ .\ 149$ & $-1''6111$ & $-0''01096$ & $42''18455$\\
$149\ .\ .\ .\ .\ 279$ & $-0.6624$ & $-0.00510$ & $42.17144$\\
$279\ .\ .\ .\ .\ 420$ & $-0.3285$ & $-0.00233$ & $42.17646$
\end{tabular}
\end{center}

Die Abnahme des Schwingungsbogens ist auch abgesehen von ihrem Zusammenhang mit der genaueren Berechnung der Schwingungsdauer noch in mehreren anderen Hinsichten von Interesse. Zunächst muss man die äußeren Umstände unterscheiden, unter denen die Nadel schwingt.

\SourceLeaf{391}Ist der Apparat zweckmäßig eingerichtet und in vollkommen gutem Zustand\HistoricalNote{391}, und befindet sich in seiner unmittelbaren Umgebung nichts, was eine beträchtliche Dämpfung der Schwingungsbewegung bewirken muss, so nehmen die Schwingungsbögen immer sehr langsam ab. Die Abnahme ist insofern regelmäßig, als sie wenigstens eine mäßig lange Zeit hindurch in geometrischer Folge verläuft, das logarithmische Dekrement also annähernd konstant ist. In unserem Beispiel ändert sich das logarithmische Dekrement während fast sechs Stunden nur von $0.00117$ bis $0.00130$, oder die Abnahme des Bogens von einer Schwingung zur nächsten schwankte zwischen $\tfrac1{371}$ und $\tfrac1{334}$. Doch zeigt die Erfahrung, dass sehr häufig sehr verschiedene Werte des logarithmischen Dekrements vorkommen. Am selben Apparat steigt es mitunter bis gegen $0.00300$ und sinkt zu anderen Zeiten auf $0.00030$, in seltenen Fällen sogar noch tiefer. Nach unseren Erfahrungen treten die Veränderungen jedoch stets nur allmählich ein. Man wird diese Abnahme also kaum allein dem Luftwiderstand zuschreiben dürfen. Die eigentliche Ursache für diese Unterschiede hat sich aber bisher unseren vielfach wiederholten Versuchen entzogen. Wir wünschen daher sehr, dass auch Beobachter an anderen Orten ihre besondere Aufmerksamkeit auf dieses noch immer rätselhafte Phänomen richten. Einen Umstand haben wir bei diesen Versuchen so oft bemerkt, dass wir ihn kaum noch für zufällig halten können, auch wenn ein ursächlicher Zusammenhang ganz unerklärlich bleibt. Die \textit{sehr} kleinen Werte des logarithmischen Dekrements treten immer nur bei bedecktem Himmel ein, die sehr großen dagegen gewöhnlich nur bei heiterem Himmel. Der Vollständigkeit halber sei noch bemerkt, dass der Apparat an einem Metalldraht und nicht an einem Seidenfaden aufgehängt ist und dass diese Versuche immer nur an windstillen Tagen angestellt wurden, folglich in diesem beiden Beziehungen sowohl ein hygrometrischer Einfluss als auch Luftzug ganz außer Frage stehen.

Ist die Nadel dagegen von einem Multiplikator umgeben, der Teil eines geschlossenen Stromkreises ist, so kommt eine weitere Ursache für die Abnahme der Schwingungsbögen hinzu. Die Bewegung der Nadel induziert nämlich in diesem Stromkreis einen galvanischen Strom, dessen Stärke am größten ist, wenn die
\expandafter\def\csname HistoricalNoteText393\endcsname{Welche Mittel man auch dazu einsetzt, man kann doch nicht damit rechnen, im strengen Sinn \textit{reine} Schwingungen um eine vertikale Achse zu erzeugen. Schon deshalb darf man von den Ergebnissen keine Übereinstimmung bis ins Kleinste erwarten, worauf es hier aber auch gar nicht ankommt.}
\expandafter\def\csname HistoricalNoteText392\endcsname{Es ist nicht bei allen Versuchen ganz gleich, da die Wirkung des Multiplikators sich mit einer Größe verbindet, die, wie oben bemerkt, selbst nicht unveränderlich ist. Die Wirkung des Multiplikators selbst hängt zudem von der Leitfähigkeit des Drahts ab, die sich mit der Temperatur ändert.}
\SourceLeaf{392}Schwingungsgeschwindigkeit der Nadel am größten ist und der die entgegengesetzte Richtung annimmt, sobald die Nadel umkehrt. Die Rückwirkung dieses Stroms auf die Nadel besteht aber immer in einer Verminderung ihrer Schwingungsgeschwindigkeit. Aus der Theorie ergibt sich, dass auch dadurch der Schwingungsbogen in nahezu genau geometrischer Folge abnehmen muss. Da aber die Stärke des induzierten Stroms auch durch den Widerstand bestimmt wird, den der \textit{ganze} Stromkreis bietet, ist die vom Multiplikator bewirkte Vergrößerung des logarithmischen Dekrements am stärksten, wenn der Stromkreis unmittelbar hinter ihm geschlossen ist. Sie fällt umso kleiner aus, je größer der hinzukommende Teil des Stromkreises ist. Bei offenem Stromkreis gibt es überhaupt keinen Einfluss, und das logarithmische Dekrement ist dasselbe, als wäre der Multiplikator ganz entfernt. Der Apparat, an dem die obigen Beobachtungen angestellt wurden, hat einen Multiplikator mit 610 Windungen. Wird er ohne weiteren Zusatz geschlossen, so steigt das logarithmische Dekrement auf etwa $0.02400$\HistoricalNote{392}. Der Schwingungsbogen sinkt dadurch schon in etwa 9 Minuten auf die Hälfte, während bei offenem Stromkreis und einem Wert des logarithmischen Dekrements, wie ihn die obigen Beobachtungen ergaben, etwa drei Stunden dafür nötig sind. Bei der vierpfündigen Nadel des magnetischen Observatorium ergibt das Schließen eines Multiplikators mit 536 Windungen ein logarithmisches Dekrement von annähernd gleicher Größe. Da jene Nadel aber eine Schwingungsdauer von $21''6$ hat, sinkt der Schwingungsbogen hier schon nach $4\tfrac12$ Minuten auf die Hälfte. Bei so starker Dämpfung können, wenn die Nadel einmal zur Ruhe gekommen ist, gar keine erhebliche Schwingungen aufkommen, sofern keine außergewöhnlichen Störungen von außen oder ungewöhnlich starke, schnelle Deklinationsänderungen eintreten. Ein so kräftiger Multiplikator hat daher neben seinen zahllosen anderen Anwendungen auch den wichtigen Nutzen, die Terminsbeobachtungen außerordentlich zu erleichtern und alle anderen Mittel zur Beruhigung der Nadel entbehrlich zu machen.

In noch viel höherem Maß leistet dies die oben [S. 372 d. B.] unter dem Namen eines Dämpfers erwähnte Vorrichtung. Der für das Magnetome\SourceLeaf{393}ter des magnetischen Observatoriums angefertigte Dämpfer besteht aus zwei länglichen rechteckigen Kupferrahmen, die je 13 Pfund wiegen. Sie lassen sich wie ein Futter in die Holzrahmen der beiden Multiplikatorhälften einschieben. Da diese Vorrichtung in sehr vielen Fällen überaus nützlich ist und auch schon mehrere von Herrn \textsc{Meyerstein} an auswärtige Beobachter gelieferte Magnetometer mit einem ähnlichen Hilfsapparat ausgestattet wurden, sind einige Bemerkungen dazu hier am Platz.

Wer einen solchen Apparat benutzen will, wird zunächst wissen wollen, wie stark er dämpft, und zwar in Zahlen. Man versetzt die Nadel zuerst in sehr große Schwingungen\HistoricalNote{393} und zeichnet die Elongationen auf, sobald sie innerhalb der Skala liegen und solange der Schwingungsbogen noch beträchtlich ist. Sind $g$, $g'$, $g''$, $g'''$ und so weiter die so erhaltenen Amplituden in Skalenteilen, so erhält man aus den Differenzen ihrer Logarithmen $\log g-\log g'$, $\log g'-\log g''$, $\log g''-\log g'''$ und so weiter ebenso viele Bestimmungen des logarithmischen Dekrements. Ändert sich die Deklination merklich, so lässt sich diese Änderung unter der Voraussetzung, dass sie während der Beobachtungen gleichförmig verläuft, durch die Formeln
\[
\log(g+g')-\log(g'+g''),\quad \log(g'+g'')-\log(g''+g''')\text{ usw.}
\]
eliminieren. Das läuft, wie man leicht sieht, auf dasselbe hinaus wie das oben [S. 383 d. B.] angegebene Verfahren.

Einer der am 9. Januar 1838 angestellten Versuche ergab beispielsweise
\begin{center}
\begin{tabular}{r|r|r|r}
Elongationen&Amplituden&Logarithmen&Differenzen\\
$134.0$&&&\\[-3pt]
&$1057.7$&$3.02436$&\\[-3pt]
$1191.7$&&&$0.31137$\\[-3pt]
&$516.4$&$2.71299$&\\[-3pt]
$675.3$&&&$0.30850$\\[-3pt]
&$253.8$&$2.40449$&\\[-3pt]
$929.1$&&&$0.31141$\\[-3pt]
&$123.9$&$2.09307$&\\[-3pt]
$805.2$&&&
\end{tabular}
\end{center}
\SourceLeaf{394}Das logarithmische Dekrement war also im Mittel $=0.31043$. Die anderen Formeln ergeben
\[
\log\frac{1574.1}{770.2}=0.31043,\quad \log\frac{770.2}{377.7}=0.30945
\]
und damit im Mittel $0.30994$. Bei diesen Versuchen war der Multiplikator nicht geschlossen. Es wirkte also nur der Dämpfer. Bei anderen Versuchen am selben Tag, bei denen der Multiplikator zugleich geschlossen war, ergab sich das logarithmische Dekrement $=0.33570$.

Da Kupfer, wenn es nicht ganz rein ist, einen wenn auch nur sehr schwachen direkten magnetischen Einfluss ausüben kann, empfiehlt es sich, den Dämpfer bei Beobachtungen zur Bestimmung der absoluten Deklination erst dann einzusetzen, wenn man sich davon überzeugt hat, dass kein merklicher Einfluss dieser Art vorliegt. Dazu beobachtet man die Stellung der Nadel abwechselnd mit und ohne Dämpfer und zugleich die Stellung einer zweiten Nadel in angemessener Entfernung, um die während der Beobachtungen auftretenden Deklinationsänderungen berücksichtigen zu können. Steht kein zweites Magnetometer zur Verfügung, kann man diese Änderungen weniger zuverlässig eliminieren, indem man die abwechselnden Bestimmungen in annähernd gleichen Zeitabständen vornimmt und jede Stellung ohne Dämpfer mit dem Mittel der vorangehenden und der folgenden Stellung mit Dämpfer vergleicht, und umgekehrt. Bei der Anfertigung des Dämpfers für das hiesige magnetische Observatorium hat Herr \textsc{Meyerstein} vorsichtshalber nur ganz neue Feilen benutzt. Der fertige Dämpfer wurde anschließend eine beträchtliche Zeit lang in verdünnte Salzsäure gelegt, dann in Lauge und schließlich mit Wasser abgespült. Versuche der beschriebenen Art haben keinen merklichen Einfluss dieses Dämpfers auf die Stellung der Nadel erkennen lassen.

Was bei stets gleich bleibendem Schwingungsbogen streng gelten würde, erfährt unter den tatsächlichen Bedingungen, bei denen der Schwingungsbogen fortwährend abnimmt, mehrere Änderungen. Sie verdienen hier eine etwas genauere Betrachtung, und sei es nur, um genau beurteilen zu können, unter welchen Umständen man sie als unmerklich ansehen darf.

Eine Abnahme des Schwingungsbogens in geometrischer Folge setzt voraus, dass der Bewegung in jedem Augenblick eine Kraft entgegenwirkt, die in einfachem Verhältnis zu
\expandafter\def\csname HistoricalNoteText395\endcsname{\textit{Streng} genommen gilt beides nur für unendlich kleine Schwingungen.}%
\SourceLeaf{395}ihrer Geschwindigkeit steht\HistoricalNote{395}. Die allgemeine Gleichung für die Schwingungsbewegung hat daher, wenn wir die Größen dritter Ordnung bezüglich des Schwingungsbogens vernachlässigen, die Form
\[
0=\frac{\mathrm{d}\mathrm{d}x}{\mathrm{d}t^2}+nn(x-p)+2\epsilon.\frac{\mathrm{d}x}{\mathrm{d}t}
\]
Dabei bezeichnen $x$ den den Stand der Nadel zur Zeit $t$ angebenden und $p$ den der Ruhestellung entsprechenden Skalenteil. $nn$ und $2\epsilon$ bezeichnen dagegen die magnetische Richtkraft und jene verzögernde Kraft, beide geteilt durch das Trägheitsmoment der Nadel. Das vollständige Integral dieser Gleichung ist
\[
x=p+Ae^{-\epsilon t}\sin\BareRoot{(nn-\epsilon\epsilon)}.(t-B)
\]
Dabei ist $e$ die Basis der natürlichen Logarithmen, und $A$ und $B$ sind die beiden durch die Integration eingeführten willkürlichen Konstanten. Ohne die verzögernde Kraft wäre das Integral
\[
x=p+A\sin n(t-B)
\]
Die Nadel schwingt also in jenem wie in diesem Fall periodisch um den Punkt $p$, doch gibt es zwei Unterschiede. Zum einen ist im zweiten Fall die größte Abweichung von der Mitte oder die halbe Amplitude konstant $=A$, während sie im ersten in geometrischer Folge abnimmt. Zum anderen wächst das Argument der periodischen Funktion im ersten Fall langsamer als im zweiten. Setzt man die Schwingungsdauer im zweiten Fall, in dem sie allein von der magnetischen Richtkraft abhängt, $=T$, im ersten $=T'$, so hat man, wenn $\pi$ die übliche Bedeutung hat,
\[
nT=\pi,\quad \BareRoot{(nn-\epsilon\epsilon)}.T'=\pi
\]
Schreibt man also der Kürze halber $n'$ statt $\BareRoot{(nn-\epsilon\epsilon)}$ und führt einen Hilfswinkel $\varphi$ ein, sodass
\[
\sin\varphi=\frac{\epsilon}{n},\quad \cos\varphi=\frac{n'}{n},\quad \operatorname{tang}\varphi=\frac{\epsilon}{n'}
\]
gilt, und bezeichnet man ferner wie oben mit $\lambda$ das logarithmische Dekrement und mit
\SourceLeaf{396}$m$ den Modul des Systems, so erhält man
\[
\begin{gathered}
T'=\frac{T}{\cos\varphi}\\
\frac{\lambda}{m}=\epsilon T'=n'\operatorname{tang}\varphi.T'=\pi\operatorname{tang}\varphi
\end{gathered}
\]
und folglich
\[
\operatorname{tang}\varphi=\frac{\lambda}{m\pi}=\frac{\lambda}{1.364376}
\]
Für $\lambda=0.02400$ und $T=42''18$ ergibt sich nach diesen Formeln $\varphi=1^0\,0'\,28''$ und $T'=42''18653$. Das bloße Schließen des Multiplikators bewirkt also nur eine geringe Verlängerung der Schwingungsdauer. Setzt man $T=20''60$, so liefern dagegen die oben bei der Verwendung des Dämpfers allein oder zusammen mit dem Multiplikator gefundenen Zahlen
\begin{center}
\begin{tabular}{r|r|r}
$\lambda=0.30994$&$\varphi=12^0\,47'\,54''$&$T'=21''12484$\\
$\lambda=0.33570$&$\varphi=13\quad49\quad22$&$T'=21.21439$
\end{tabular}
\end{center}

Die Beobachtungen stimmen mit dieser berechneten Verlängerung der Schwingungsdauer so genau überein, wie man es von der kleinen Zahl von Schwingungen, auf die man sich dabei beschränken muss, nur erwarten kann.

Bei abnehmenden Schwingungsbögen stimmen die wahren Zeiten der Elongationen mit den aus entsprechenden Stellungen abgeleiteten Zeiten nicht genau überein. Bei so großen logarithmischen Dekrementen, wie sie bei der Verwendung eines Dämpfers vorkommen, wird dieser Unterschied ziemlich beträchtlich.

Im oben angegebenen Integral bezeichnet $B$ offenbar den Zeitpunkt eines Durchgangs durch die Ruhestellung $p$. Da es gleichgültig ist, von welchem Augenblick an die Zeit gezählt wird, wollen wir der Einfachheit halber $B=0$ setzen. Unsere Formel lautet dann
\[
x=p+Ae^{-\epsilon t}\sin n't
\]
Der nächste Durchgang durch $p$ nach demjenigen bei $t=0$ findet bei $n't=180^0$ oder $t=T'$ statt. Die aus diesen entsprechenden Beobachtungen unmittelbar abgeleitete Elongationszeit ist also $t=\tfrac12T'$, während der wirkliche Stillstand schon früher eintritt. Für $\frac{\mathrm{d}x}{\mathrm{d}t}=0$ hat man nämlich
\[
0=Ae^{-\epsilon t}(-\epsilon\sin n't+n'\cos n't)
\]
\SourceLeaf{397}Mithin
\[
\operatorname{cotang}n't=\frac{\epsilon}{n'}=\operatorname{tang}\varphi
\]
Daher ist der erste positive Wert $n't=\tfrac12\pi-\varphi$, und es gilt $t=\tfrac12T'-\frac{\varphi T'}{\pi}$ oder, wenn $\varphi$ in Grad ausgedrückt ist,
\[
t=\tfrac12T'-\frac{\varphi}{180^0}.T'
\]

Offenbar ergibt sich dieselbe Differenz bei der nächsten Elongation. Aus den oben angegebenen Zahlen ergibt sie sich $=0''23$ für den fünfundzwanzigpfündigen Stab bei Verwendung des Multiplikators, $=1''50$ für das Magnetometer des M. O., wenn der Dämpfer allein benutzt wird, und $=1''63$, wenn Dämpfer und Multiplikator zugleich verwendet werden.

Die wirklichen Stillstände treten um eine konstante, von der Größe des Schwingungsbogens unabhängige Zeit früher ein als die aus aufeinanderfolgenden Durchgängen durch die Ruhestellung $p$ erschlossenen Zeitpunkte. Man kann diese Zeitpunkte also ohne Weiteres beibehalten, da es bei ihrer Verwendung zur Bestimmung der Schwingungsdauer nur auf die \textit{Unterschiede} der Elongationszeiten ankommt. Man muss nur darauf achten, den Punkt $p$ selbst oder einen sehr nahe gelegenen Punkt zum Beobachten zu wählen. Kleine Schwingungsbögen sind umso mehr auszuschließen, als bei ihnen im Fall großer logarithmischer Dekremente schon eine geringe Abweichung vom richtigen $p$ einen merklichen Fehler verursachen würde. Obwohl diese Bedingung leicht zu erfüllen ist, mag hier doch noch die allgemeine Formel für den Fehler der aus entsprechenden Durchgängen erschlossenen Elongationszeiten Platz finden.

Es sei $u$ die halbe Zeitspanne zwischen zwei entsprechenden Durchgängen durch einen Punkt $x$ und $\tfrac12T'-\delta$ das Mittel der Durchgangszeiten oder die daraus erschlossene Elongationszeit. Die Durchgangszeiten selbst sind also $\tfrac12T'-u-\delta$ und $\tfrac12T'+u-\delta$, und wir haben folglich
\[
\begin{aligned}
x&=p+Ae^{-\epsilon(\tfrac12T'-u-\delta)}\sin n'(\tfrac12T'-u-\delta)\\
x&=p+Ae^{-\epsilon(\tfrac12T'+u-\delta)}\sin n'(\tfrac12T'+u-\delta)
\end{aligned}
\]
woraus wegen $\tfrac12n'T'=\tfrac12\pi$ folgt
\[
e^{2\epsilon u}\cos n'(u+\delta)=\cos n'(u-\delta)
\]
\SourceLeaf{398}und somit
\[
\operatorname{tang}n'\delta=\frac{e^{2\epsilon u}-1}{(e^{2\epsilon u}+1)\operatorname{tang}n'u}
\]

Für $u=\tfrac12T'$ ergibt diese Formel $\delta=0$. Das ist der Fall, in dem die Ruhestellung $p$ selbst für die Durchgänge gewählt ist. Die Annahme eines unendlich kleinen Werts entspricht dagegen dem wirklichen Stillstandspunkt. Die Formel ergibt hier
\[
\operatorname{tang}n'\delta=\frac{\epsilon}{n'}=\operatorname{tang}\varphi
\]
Also $\delta=\frac{\varphi}{n'}=\frac{\varphi T'}{\pi}$, in Übereinstimmung mit dem oben gefundenen Ergebnis.

Schließlich muss man bei abnehmenden Schwingungsbögen auch die Berechnung der Beobachtungen zur Ruhestellung der Nadel ändern. Das macht sich freilich erst bei einer sehr starken Abnahme der Schwingungen bemerkbar.

Die Nadelstellungen $x$, $x'$, die zwei um eine Schwingungsdauer auseinanderliegenden Zeiten $t$, $t+T'$ entsprechen, haben die Werte
\[
\begin{aligned}
x&=p+Ae^{-\epsilon t}\sin n't\\
x'&=p+Ae^{-\epsilon t-\epsilon T'}\sin(n't+n'T')
\end{aligned}
\]
oder, weil $n'T'=\pi$,
\[
x'=p-Ae^{-\epsilon t-\epsilon T'}\sin n't
\]
oder, wenn wir wie oben mit $\theta$ den Bruch bezeichnen, dessen Briggsscher Logarithmus $-\lambda$ und dessen natürlicher Logarithmus daher $\epsilon T'$ ist,
\[
x'=p-\theta Ae^{-\epsilon t}\sin n't
\]

Um $p$ zu finden, darf man also nicht mehr das arithmetische Mittel zwischen $x$ und $x'$ nehmen, sondern muss die Differenz zwischen $x$ und $x$ im Verhältnis von 1 zu $\theta$ aufteilen. Man erhält also
\[
p=\frac{\theta x+x'}{1+\theta}=x+\frac{1}{1+\theta}.(x'-x)=x'-\frac{\theta}{1+\theta}.(x'-x)
\]
Da bei diesen Beobachtungen die Differenz $x'-x$ übrigens immer sehr klein ist, kann man sich die Rechnung erleichtern, indem man statt $\frac{\theta}{1+\theta}$ einen nahe liegenden, durch kleine Zahlen ausdrückbaren Bruch
\SourceLeaf{399}verwendet. So kann man etwa für $\lambda=0.30994$ statt des genaueren Werts $0.3288$ den Näherungswert $\tfrac13$ wählen.

Dabei stellt sich nun die Frage, für welchen Augenblick dieses Ergebnis als gültig anzusehen ist. Wenn die Schwingungsbögen nur sehr langsam abnehmen, nimmt man an, dass das einfache Mittel der Skalenteile dem einfachen Mittel der Zeiten entspricht. Entsprechend könnte es scheinen, als sei bei ungleich aufgeteilter Differenz der Stellungen die Zwischenzeit im selben Verhältnis zu teilen, also die Stellung $x+\frac{x'-x}{1+\theta}$ als für $t+\frac{T'}{1+\theta}$ gültig anzusehen. Das ist aber theoretisch nicht richtig, und eine genauere Erörterung, die in gewöhnlichen Fällen praktisch ganz unerheblich ist, verdient hier vielleicht doch aus theoretischen Gründen ihren Platz.

Offenbar kommt der Gültigkeitszeitpunkt nur insofern in Betracht, als man die magnetische Deklination in der Zwischenzeit zwischen den beiden Aufzeichnungen nicht als konstant ansehen will. Man kann aber immer annehmen, dass sie sich während dieser Zeit gleichförmig ändert, und muss dies auch annehmen, wenn die Rechnung eine bestimmte Grundlage erhalten soll\quad In diesem Fall hat unsere Grundgleichung also die Form
\[
0=\frac{\mathrm{d}\mathrm{d}x}{\mathrm{d}t^2}+nn(x-p-\alpha t)+2\epsilon.\frac{\mathrm{d}x}{\mathrm{d}t}
\]
und ihr vollständiges Integral ist
\[
x=p-\frac{2\alpha\epsilon}{nn}+\alpha t+Ae^{-\epsilon t}\sin n'(t-B)
\]
wenn wie oben $n'$ für $\BareRoot{(nn-\epsilon\epsilon)}$ gesetzt wird. Bezeichnet also $x'$ die Stellung zur Zeit $t+T'$, so erhält man mit $\theta$ in seiner bisherigen Bedeutung
\[
x'=p-\frac{2\alpha\epsilon}{nn}+\alpha t+\alpha T'-\theta Ae^{-\epsilon t}\sin n'(t-B)
\]
und folglich
\[
\frac{\theta x+x'}{1+\theta}=p-\frac{2\alpha\epsilon}{nn}+\alpha t+\frac{\alpha T'}{1+\theta}
\]
Dieses Ergebnis ist demnach die Ruhestellung für die Zeit
\[
t+\frac{T'}{1+\theta}-\frac{2\epsilon}{nn}=t+\frac{T'}{1+\theta}-\frac{\sin2\varphi.T'}{\pi}
\]
Für $\lambda=0.30994$ gilt also das nach der obigen Vorschrift berechnete Ergebnis zur Zeit $=t+0.5337T'$, für $\lambda=0.33570$ dagegen zur Zeit $=t+0.5365T'$.
\SourceLeaf{400}Man sieht also, dass selbst bei so starker Dämpfung der Gültigkeitszeitpunkt nur wenig vom einfachen Mittel der Zeiten abweicht.

Alles bisher Entwickelte beruht auf der Voraussetzung, dass $\epsilon$ kleiner als $n$ ist. Im entgegengesetzten Fall nimmt das Integral der Grundgleichung eine andere Form an. Anstelle des Glieds $Ae^{-\epsilon t}\sin\BareRoot{(nn-\epsilon\epsilon)}.(t-B)$ erhält man, wenn $\epsilon$ größer als $n$ ist, zwei Glieder der Form
\[
Ae^{-(\epsilon+\BareRoot{(\epsilon\epsilon-nn)})t}+Be^{-(\epsilon-\BareRoot{(\epsilon\epsilon-nn)})t}
\]
und wenn $\epsilon=n$ ist, ein Glied der Form
\[
(A+Bt)e^{-\epsilon t}
\]

In beiden Fällen ist die Bewegung also nicht mehr periodisch, und die Stellung nähert sich asymptotisch der Ruhestellung. Für unseren Dämpfer ist $\frac{\epsilon}{n}=0.22152$. Man müsste also einen mehr als $4\tfrac12$-mal so stark wirkenden Dämpfer verwenden, um diese Wirkung zu erzielen. Dazu würde es offenbar nicht genügen, die Metallmenge im selben Verhältnis zu vergrößern, da das zusätzliche Metall außen angebracht werden müsste und die äußeren Schichten des Metallrahmens verhältnismäßig weniger zur Induktionswirkung beitragen als die inneren. Es wäre aber auch nicht ratsam, eine so starke Dämpfung anzuwenden, dass die Bewegung aufhört, periodisch zu sein. Zum einen nähert sich die Stellung der Ruhestellung wieder langsamer, sobald $\epsilon$ den Grenzwert $n$ überschreitet. Zum anderen ginge der wesentliche Vorteil verloren, die Ruhestellung aus zwei beliebigen, um $T'$ auseinanderliegenden Aufzeichnungen bequem berechnen zu können.

